\documentclass[11pt]{article}

\usepackage[margin=2.5cm]{geometry}
\usepackage[utf8]{inputenc}
\usepackage[T1]{fontenc}
\usepackage{amsmath,amssymb}
\usepackage{booktabs}
\usepackage{array}
\usepackage{xcolor}
\usepackage{colortbl}
\usepackage{titlesec}
\usepackage{enumitem}
\usepackage{parskip}
\usepackage{longtable}
\usepackage{fancyhdr}
\usepackage{tcolorbox}
\usepackage{graphicx}
\usepackage[colorlinks=true,linkcolor=blue!45!black,urlcolor=blue!45!black]{hyperref}

\definecolor{crit}{RGB}{155,25,25}
\definecolor{warn}{RGB}{180,110,0}
\definecolor{good}{RGB}{20,110,60}
\definecolor{hdr}{RGB}{30,55,105}

\titleformat{\section}{\normalfont\Large\bfseries\color{hdr}}{\thesection}{0.6em}{}
\titleformat{\subsection}{\normalfont\large\bfseries\color{hdr}}{\thesubsection}{0.6em}{}

\pagestyle{fancy}
\fancyhf{}
\fancyhead[L]{\small\textit{Análisis estadístico de datos reales -- CICIDS-2018}}
\fancyhead[R]{\small\thepage}
\renewcommand{\headrulewidth}{0.4pt}

\newcommand{\code}[1]{\texttt{\small #1}}

\newtcolorbox{cbox}[2]{colback=#1!5,colframe=#1,fonttitle=\bfseries,title={#2},
  boxrule=0.9pt,left=5pt,right=5pt,top=4pt,bottom=4pt}

\title{\textbf{\Large Primer análisis estadístico de datos reales}\\[4pt]
\Large \textit{Documentación de columnas y contraste contra los supuestos del simulador}\\[8pt]
\normalsize CICIDS-2018 real $\cdot$ dos archivos $\cdot$ 8.28M flujos}
\author{Repositorio \code{Ironsss/iag-symmetry}}
\date{2026-08-29}

\begin{document}
\maketitle
\thispagestyle{empty}

\begin{cbox}{hdr}{Qué es este documento y por qué existe}
Hasta esta sesión, todo número que el proyecto citó sobre comportamiento de features,
asimetría per-flujo, o criterios de Fase 0 provenía de \code{src/simulator/agents.py}
--- tráfico sintético. Esta sesión descargó, por primera vez, dos archivos reales de
CICIDS-2018 desde el bucket público. Este documento mide, sobre esos datos reales,
exactamente lo que el resto del repositorio había medido sólo en simulación, y reporta
la brecha donde la hay. La brecha es grande: el hallazgo central (\S\ref{sec:contraste})
es que la asimetría per-flujo, que en el simulador separaba ataques con AUC 0.59--0.95,
cae a un rango mucho más débil en tráfico real.

\medskip
Ampliado con tres cosas que ninguna tabla de un solo número mostraba: el grafo de
interacciones real (\S\ref{sec:grafo} --- nunca se había visualizado, pese a ser el
objeto central del framing original del proyecto), clustering no supervisado sobre
cúmulos de flujos, no flujo por flujo (\S\ref{sec:clustering}), y el primer benchmark
supervisado real (\S\ref{sec:benchmark} --- los cinco modelos de
\code{src/models/benchmark.py} nunca habían corrido contra datos reales), y el experimento
preregistrado H2 bajo restricción de información, corrido por primera vez contra datos reales
(\S\ref{sec:h2} --- el gate de potencia se cumple por primera vez en la historia del proyecto,
aunque el de saturación todavía no cede), y los 10 archivos del bucket, no sólo 2
(\S\ref{sec:completitud} --- el AUC per-flujo invertido resulta ser específico de ataques de
\emph{flood}, no un fenómeno general). Cierra con una síntesis que conecta desbalanceo,
asimetría, topología, y estado del arte (\S\ref{sec:sintesis}).

\medskip
Cada número de este documento es reproducible con \code{python scripts/analyze\_real\_data.py
data/raw/} y \code{python scripts/explore\_real\_data\_structure.py data/raw/Thuesday-20-02-2018\_TrafficForML\_CICFlowMeter.csv},
cuya salida completa queda en \code{results/metrics/real\_data\_summary.json} y
\code{results/metrics/real\_data\_clusters.json}.
\end{cbox}

\tableofcontents
\newpage

% ============================================================
\section{Procedencia y alcance}
% ============================================================

Dos archivos, ambos descargados en esta sesión del bucket público
\code{s3://cse-cic-ids2018/} (\S5 de \code{docs/onboarding/setup/SETUP-LOCAL.md}):

\begin{center}
\begin{tabular}{@{}lll@{}}
\toprule
& \textbf{Thursday-01-03-2018} & \textbf{Thuesday-20-02-2018} \\
\midrule
Día / ataque & Infiltración & DDoS-LOIC-HTTP \\
Tamaño & 108 MB & 3.8 GB \\
Filas & 331.126 & 7.948.746 \\
Esquema & sin \code{Src IP}/\code{Dst IP} & ancho (con IPs) \\
\bottomrule
\end{tabular}
\end{center}

\code{Thuesday-20-02-2018} es, verificado contra los 10 archivos del bucket vía \code{GET}
por rango HTTP (sin descargarlos completos), el \textbf{único} de los diez con el esquema
ancho que el grafo necesita. No hay un segundo día con IPs para concatenar si un criterio
de ventana no se cumple con éste solo --- eso ya se documentó en
\code{docs/onboarding/setup/SETUP-LOCAL.md} y es relevante otra vez en \S\ref{sec:phase0}.

% ============================================================
\section{Documentación de columnas}
% ============================================================

\code{Thuesday-20-02-2018} trae 84 columnas crudas; \code{Thursday-01-03-2018} trae 80
(le faltan \code{Flow ID}, \code{Src IP}, \code{Src Port}, \code{Dst IP}). De esas, el
repositorio ya distingue tres roles en \code{src/ingest/loader.py}:

\begin{center}
\begin{tabular}{@{}lp{9.5cm}r@{}}
\toprule
\textbf{Rol} & \textbf{Qué es} & \textbf{Cuántas} \\
\midrule
identificador/IP & \code{Flow ID}, \code{Src IP}, \code{Src Port}, \code{Dst IP} & 4 (solo en el archivo DDoS) \\
timestamp & \code{Timestamp} & 1 \\
label & \code{Label} & 1 \\
\code{CLASSIC\_FEATURES} & el subconjunto curado que \code{src/models/benchmark.py} entrena & 49 \\
otras mapeadas & normalizadas por \code{COLUMN\_RENAME} pero fuera del baseline curado (p.\,ej.\ \code{subflow\_*}, \code{init\_*}, \code{cwe\_flag\_cnt}) & 29 \\
\bottomrule
\end{tabular}
\end{center}

De las 49 \code{CLASSIC\_FEATURES}, \textbf{24 son direccionales} (\code{"fwd"} o \code{"bwd"}
en el nombre) --- el mismo 24/49 que \code{scripts/validate\_claims.py} ya verificaba
estructuralmente, ahora confirmado además contra los datos reales: las 49 están presentes y
con datos válidos en ambos archivos, no sólo en el nombre de la columna.

% ============================================================
\section{Calidad de datos, medida}
% ============================================================

\begin{cbox}{crit}{Hallazgo: headers duplicados corrompían el tipo de dato en el loader de producción}
\code{Thursday-01-03-2018} es la concatenación de $\sim$26 exports de CICFlowMeter, cada
uno con su propia fila de header repetida en medio del archivo (documentado primero en
\code{docs/onboarding/setup/SETUP-LOCAL.md} durante el preflight). Lo que esta sesión
encontró después es más grave: sin una anulación de esquema explícita, la inferencia de
tipos de \code{polars} veía el texto del header como un valor no numérico dentro de lo que
debía ser una columna numérica, y --- en vez de anular sólo ese valor --- \textbf{terminaba infiriendo
la columna entera como \code{String}}. Le pasaba a las 49 \code{CLASSIC\_FEATURES} en este
archivo. \code{\_clean\_data()} sólo limpia columnas de tipo numérico, así que su limpieza
de infinitos/NaN nunca corría sobre este archivo. Corregido en \code{src/ingest/loader.py}
(PR \#17): se lee el header real como texto plano y se fuerza \code{Float64} vía
\code{schema\_overrides}, de modo que \code{ignore\_errors=True} anule sólo el valor
puntual, no el tipo completo de la columna.
\end{cbox}

\begin{center}
\begin{tabular}{@{}lrr@{}}
\toprule
\textbf{Medida} & \textbf{Infiltración} & \textbf{DDoS} \\
\midrule
Filas de header incrustadas & 25 & 0 \\
\code{flow\_byts\_per\_s} = $\infty$ (antes de limpiar) & 1.085 & 22.686 \\
\code{flow\_pkts\_per\_s} = $\infty$ (antes de limpiar) & 2.919 & 59.453 \\
Filas tras deduplicar (\code{.unique()}) & 331.028 & 7.948.746 \\
Features degeneradas (\code{std}=0 o >99\% nulas) & \code{bwd\_psh\_flags} & \code{bwd\_psh\_flags} \\
\bottomrule
\end{tabular}
\end{center}

Los infinitos son reales, no un artefacto de parseo: \code{flow\_byts\_per\_s} y
\code{flow\_pkts\_per\_s} dividen bytes/paquetes por duración, y CICFlowMeter emite
\code{Infinity} literal cuando la duración es cero. \code{\_clean\_data()} los reemplaza
por 0.0 --- una aproximación documentada, no una suposición nueva de este análisis.

Queda un residuo menor sin resolver: un puñado de filas conserva un \code{null} genuino
(no infinito) en todas las \code{CLASSIC\_FEATURES} tras \code{\_clean\_data()}, artefacto
de las filas de header incrustadas que \code{ignore\_errors} anuló --- \code{\_clean\_data}
limpia infinitos/NaN, no nulos genuinos. Este análisis las excluye explícitamente
(\code{.drop\_nulls()}) antes de cualquier cómputo sensible a \code{NaN} (AUC, correlación).
No se filtran en el pipeline general; queda como mejora de bajo costo, no bloqueante.

% ============================================================
\section{Auditoría de features direccionales en datos reales}
% ============================================================

Más allá del conteo estructural 24/49, cada feature se midió por tasa de nulos y varianza
real (no sólo si el nombre existe). Sobre \code{Thursday-01-03-2018}: de 49, sólo
\textbf{\code{bwd\_psh\_flags}} resultó degenerada (constante o $>$99\% nula) --- el resto
de las 24 direccionales tienen datos genuinamente variables en este archivo. El detalle
completo, columna por columna (\code{count}, \code{null\_rate}, \code{mean}, \code{std},
\code{p50}, \code{p90}, \code{p99}, \code{max}), está en
\code{results/metrics/real\_data\_summary.json} bajo \code{numeric\_summary} --- 49 filas
por archivo, demasiado para reproducir aquí sin perder legibilidad; \S\ref{sec:corr} resume
lo que importa de esa tabla (redundancia), y las figuras muestran la forma de las
distribuciones más sesgadas.

% ============================================================
\section{Estructura de correlación real}
\label{sec:corr}
% ============================================================

Sobre una muestra sistemática determinística (cada $k$-ésima fila, sin RNG --- ver
\S\ref{sec:metodo}) de las \code{CLASSIC\_FEATURES} no degeneradas, agrupadas por
\code{src/audit/ablation.py::correlation\_groups(threshold=0.80)} --- el mismo umbral que
ese módulo ya midió como el punto donde el daño a la importancia por permutación es severo
(\S3 de \code{CLAUDE.md}):

\begin{center}
\begin{tabular}{@{}lp{10.5cm}@{}}
\toprule
\textbf{Archivo} & \textbf{Grupos redundantes ($\geq$2 features, corr $\geq 0.80$)} \\
\midrule
Infiltración & 6 grupos, incl.\ \{\code{tot\_fwd\_pkts}, \code{tot\_bwd\_pkts},
  \code{totlen\_bwd\_pkts}, \code{fwd\_header\_len}, \code{bwd\_header\_len}\} y
  \{\code{fwd\_pkt\_len\_max/std}, \code{bwd\_pkt\_len\_*}, \code{pkt\_len\_*},
  \code{pkt\_size\_avg}\} \\
DDoS & 9 grupos (muestra sistemática, $k$=10, 794.875 filas), el más grande:
  \{\code{fwd\_pkt\_len\_*}, \code{bwd\_pkt\_len\_*}, \code{pkt\_len\_*},
  \code{psh\_flag\_cnt}, \code{pkt\_size\_avg}\} (12 features) \\
\bottomrule
\end{tabular}
\end{center}

La estructura de redundancia en datos reales es consistente con lo esperado de features
de conteo/tamaño de paquete derivadas del mismo flujo (misma familia que
\code{src/audit/ablation.py} ya documentaba con datos sintéticos) --- confirma que el
umbral 0.80 sigue siendo razonable contra datos reales, no sólo contra el generador
controlado con el que se midió originalmente.

\begin{figure}[h]
\centering
\includegraphics[width=0.75\textwidth]{../../results/figures/real_data_thursday-01-03-2018_correlation.png}
\caption{Correlación entre \code{CLASSIC\_FEATURES}, Infiltración real.}
\end{figure}

\begin{figure}[h]
\centering
\includegraphics[width=0.75\textwidth]{../../results/figures/real_data_thuesday-20-02-2018_correlation.png}
\caption{Correlación entre \code{CLASSIC\_FEATURES}, DDoS-LOIC-HTTP real.}
\end{figure}

% ============================================================
\section{El contraste central: AUC per-flujo real vs.\ documentado}
\label{sec:contraste}
% ============================================================

\begin{cbox}{crit}{Esta es la comparación que motivó el documento}
\code{scripts/validate\_claims.py::claim\_perflow\_auc} mide, sobre el simulador, AUC
0.9978 (PortScan) y 0.9464 (DDoS) para la asimetría per-flujo
$\log(1+\text{fwd}) - \log(1+\text{bwd})$. \code{docs/metodologia/README.md} cita,
también del simulador, un rango 0.587--0.686 para infiltración. Misma fórmula exacta,
ahora sobre flujos reales:
\end{cbox}

\begin{center}
\begin{tabular}{@{}lrrr@{}}
\toprule
\textbf{Archivo (real)} & \textbf{AUC real} & \textbf{AUC simulador (documentado)} & \textbf{n flujos} \\
\midrule
Infiltración & 0.537 & 0.587--0.686 & 331.027 (28.1\% ataque) \\
DDoS-LOIC-HTTP & \textbf{0.381} & 0.9464 & 7.948.746 (7.25\% ataque) \\
\bottomrule
\end{tabular}
\end{center}

Infiltración cae apenas por encima del azar (0.5) --- ya es más débil que el rango
documentado en el simulador. Pero DDoS es el resultado que no se podía anticipar
razonando: \textbf{0.381, por debajo de 0.5}. No es "sin señal" --- es señal
\emph{sistemáticamente invertida}: en este tráfico real, las ventanas de ataque tienen en
promedio \emph{menos} bytes enviados relativo a recibidos que el tráfico benigno, exactamente
lo opuesto al supuesto (asimetría alta = probable ataque) sobre el que se construyó
$A_{\text{dir}}$ y sobre el que el simulador entrena su tráfico de ataque. Una lectura
mecánica plausible: LOIC-HTTP envía peticiones pequeñas y recibe respuestas/páginas de
error --- el atacante manda poco y recibe más, al revés de un escaneo de puertos o una
exfiltración. Es una hipótesis de por qué, no una prueba causal; lo que sí es directo es
que \textbf{la dirección del efecto que el simulador asume no es universal}, y que
asumirla sin medir es exactamente el error que \code{CLAUDE.md} ya cataloga (mistake \#3:
``a default was reasoned rather than measured''). \S\ref{sec:completitud} confirma esta
hipótesis contra los otros 8 tipos de ataque del dataset: la inversión resulta específica
de ataques de \emph{flood} (DoS/DDoS), no un fenómeno general.

Ninguno de los dos números es una confirmación del simulador; ambos van en la dirección
que \code{docs/metodologia/REPLANTEO-OPERACIONAL.md} ya predecía (la per-flujo no
sobrevive fuera del régimen limpio de CICIDS), pero con una magnitud --- y, en DDoS, una
\emph{dirección} --- que el argumento teórico no anticipaba.

% ============================================================
\section{Barrido de criterios de Fase 0}
\label{sec:phase0}
% ============================================================

Extiende el chequeo único (\code{window=60s}, $\tau$=0.01) que \code{scripts/preflight\_data.py}
ya corrió sobre \code{Thuesday-20-02-2018}, a una grilla de \code{window\_seconds}
$\in \{10,30,60,120,300\}$ $\times$ $\tau \in \{0.01, 0.05, 0.1\}$, para saber si el
resultado depende del parámetro elegido o es robusto.

\begin{center}
\begin{tabular}{@{}rrrrl@{}}
\toprule
\textbf{window (s)} & \textbf{n ventanas} & \textbf{n ataque (min clase)} & \textbf{\% parcial} & \textbf{Criterios} \\
\midrule
10  & 3.873 & \textbf{383} (cumple $\geq$250) & 9.76\% & 3 sí / 4 no (9.76\%$<$10\%) \\
30  & 1.418 & 127 & 8.89\% & 3 no / 4 no \\
60  & 719   & 63  & 8.76\% & 3 no / 4 no \\
120 & 360   & 32  & 8.89\% & 3 no / 4 no \\
300 & 144   & 14  & 9.72\% & 3 no / 4 no \\
\bottomrule
\end{tabular}
\end{center}

($\tau$ no mueve la aguja dentro de este rango --- 0.01/0.05/0.1 dan resultados
prácticamente idénticos en cada fila; la tabla muestra $\tau$=0.01.) Dos cosas cambian la
lectura de \code{SETUP-LOCAL.md}: primero, achicar la ventana a \textbf{10s sí cumple el
criterio 3} (383 ventanas de ataque, sobre el mínimo de 250) --- la conclusión ``ningún
parámetro alcanza'' no era cierta, sólo lo era en el punto (60s) que el preflight probó
por defecto. Segundo, el criterio 4 queda \textbf{cerca pero no cumplido en todo el
rango} (8.76\%--9.76\% contra el 10\% requerido) --- no es un fracaso categórico, es un
margen estrecho y consistente, lo que sugiere que vale la pena probar ventanas aún más
finas (5s, 2s) antes de concluir que el criterio 4 es inalcanzable con este único
archivo.

% ============================================================
\section{Pipeline actual sobre datos reales, por primera vez}
\label{sec:ddos-results}
% ============================================================

\begin{cbox}{good}{Por qué esta sección corre \code{src/features/windows.py}, no \code{src/iag/}}
\code{CLAUDE.md} es explícito: \code{src/iag/} es \textit{``el índice retirado, mantenido
como baseline auditado, no como propuesta''}; \code{src/features/} (tiers) es
\textit{``la tesis, como código''}. La pregunta preregistrada vigente es H2
(\code{docs/metodologia/REPLANTEO-OPERACIONAL.md}): si la topología degrada menos que la
asimetría per-flujo bajo restricción. Correr el índice retirado sobre datos reales no
respondería nada que el proyecto siga preguntando; correr
\code{window\_features(tier=DIRECTIONAL)} --- que calcula topología, per-flujo, y el
control C1 de tamaño en una sola llamada, sobre datos reales por primera vez --- sí.
\end{cbox}

Sobre la configuración por defecto (60s, $\tau$=0.01: 719 ventanas, 63 de ataque),
\code{window\_features(tier=DIRECTIONAL)} da, por primera vez, el AUC a nivel de
\emph{ventana} de cada familia de features contra la misma etiqueta:

\begin{center}
\begin{tabular}{@{}llr@{}}
\toprule
\textbf{Familia} & \textbf{Feature} & \textbf{AUC (ventana, n=719, 63 positivas)} \\
\midrule
Per-flujo & \code{fl\_mean} & \textbf{0.0000} \\
Per-flujo & \code{fl\_p90} & 0.0328 \\
Per-flujo & \code{fl\_frac\_gt1} & 0.0517 \\
Per-flujo & \code{fl\_std} & 0.9926 \\
\midrule
Topología & \code{top\_indeg\_gini} & 0.7232 \\
Topología & \code{top\_edges\_per\_node} & 0.7001 \\
Topología & \code{top\_reciprocity} & 0.6096 \\
Topología & \code{top\_fanout\_max\_frac} & 0.2882 \\
Topología & \code{top\_maxindeg\_frac} & 0.2604 \\
\midrule
\rowcolor{crit!8} Control C1 (tamaño) & \code{sz\_nodes} & 0.7810 \\
\rowcolor{crit!8} Control C1 (tamaño) & \code{sz\_edges} & 0.7703 \\
\bottomrule
\end{tabular}
\end{center}

\begin{cbox}{warn}{No leer esto como H2 confirmada}
Dos advertencias, ambas obligatorias por las propias reglas del proyecto
(\code{CLAUDE.md}, mistake \#1 y \#2):

\textbf{Potencia.} 63 ventanas positivas está muy por debajo de las $\geq$250/clase que
\code{docs/metodologia/README.md} exige antes de confiar en un AUC. Estos números son
descriptivos de esta corrida puntual, no una prueba confirmatoria de H2.

\textbf{El control C1 no se despega de la topología.} \code{sz\_nodes} (0.7810) y
\code{sz\_edges} (0.7703) --- el control de ``sólo tamaño'', que nunca debería ser
tratamiento --- puntúan \emph{tan alto o más} que \code{top\_indeg\_gini} (0.7232) y
\code{top\_edges\_per\_node} (0.7001). Es exactamente el patrón que produjo el confound de
AUC 1.0000 documentado en \code{docs/metodologia/README.md} (positivos y negativos con
\code{n\_nodes} sin solapamiento). No se puede afirmar que la topología aporte algo que el
tamaño de la ventana no explique ya, con esta corrida sola.

Lo único que se sostiene sin reserva es la asimetría del resultado per-flujo:
\code{fl\_mean} en \textbf{0.0000} (inversión casi perfecta, consistente con el AUC a nivel
de flujo de \S\ref{sec:contraste}) contra \code{fl\_std} en 0.9926 --- la \emph{dispersión}
del ratio por ventana separa casi perfectamente, aunque su \emph{media} esté invertida.
Sugiere que la ventana de ataque no cambia el promedio de la asimetría, sino cuánto varía
--- una pista concreta para la siguiente iteración de H2, no una respuesta.
\end{cbox}

% ============================================================
\section{El grafo de interacciones real}
\label{sec:grafo}
% ============================================================

Todo el framing original de este proyecto (IAG = índice de asimetría de \emph{grafos}) se
construyó y se retiró sin que nadie viera nunca un grafo real. Esta sección lo muestra:
nodos = hosts, aristas = flujos agregados por par \code{(src, dst)}, dibujados con
\code{networkx} (reintroducido como dependencia --- se había retirado el 2026-08-25 por no
tener uso real; ahora lo tiene) sobre dos ventanas de 60s reales del archivo DDoS: una con
tráfico de ataque, una completamente benigna.

\begin{center}
\begin{tabular}{@{}lrrr@{}}
\toprule
\textbf{Ventana} & \textbf{hosts} & \textbf{aristas} & \textbf{aristas con ataque} \\
\midrule
Ataque  & 1.579 & 5.342 & \textbf{10} (0.19\%) \\
Benigna & 1.733 & 5.649 & 0 \\
\bottomrule
\end{tabular}
\end{center}

\begin{figure}[ht]
\centering
\includegraphics[width=0.85\textwidth]{../../results/figures/real_data_interaction_graph_attack.png}
\caption{Ventana de ataque real (100 hosts de mayor grado + todos los adyacentes a una
arista de ataque, de 1.579 totales). El componente rojo, aislado, es el ataque: atacantes
hablándole a un único destino, sin ninguna conexión con el resto del tráfico de fondo.}
\end{figure}

\begin{figure}[ht]
\centering
\includegraphics[width=0.85\textwidth]{../../results/figures/real_data_interaction_graph_benign.png}
\caption{Ventana benigna real (100 hosts de mayor grado, de 1.733 totales). Estructura de
hub-and-spoke: un puñado de servidores concentran la mayoría de las conexiones.}
\end{figure}

\begin{cbox}{good}{Lo que el grafo muestra que ninguna tabla mostraba}
Sólo \textbf{10 de 5.342 aristas (0.19\%)} en la ventana de ataque llevan tráfico de
ataque, y ese componente queda \textbf{topológicamente desconectado} del resto de la
red --- ni un solo host benigno de fondo interactúa con el par atacante/víctima en esta
ventana. Es la misma dilución que \code{docs/review/} ya documentó de forma numérica
(\code{A\_dir} diluido por tráfico de fondo): acá se ve. También explica, en parte, por
qué el control C1 (tamaño) puntúa tan alto en \S\ref{sec:ddos-results} --- una ventana de
ataque con un componente extra, aislado, tiene por construcción más estructura que
contar que una puramente benigna, sin que eso sea la asimetría relacional que el proyecto
busca medir.
\end{cbox}

% ============================================================
\section{Clustering: pensar en cúmulos de filas, no en filas}
\label{sec:clustering}
% ============================================================

Todo lo anterior mide una fila (un flujo) o una ventana a la vez. Esta sección agrupa: sobre
la misma muestra sistemática determinística de \S\ref{sec:corr} (305.721 filas, cada
26$^{\text a}$ fila), \code{KMeans} sobre las 48 \code{CLASSIC\_FEATURES} no degeneradas,
con $k$ elegido por silhouette score ($k \in [2,8]$, sobre un sub-muestreo de 5.000 filas
por costo $O(n^2)$ de silhouette).

\begin{cbox}{warn}{El preprocesamiento importó más de lo esperado}
Dos intentos previos de proyectar esto a 2D con PCA fallaron de forma instructiva:
\code{StandardScaler} solo dejó que 1-2 filas con un \code{flow\_duration}/\code{flow\_byts\_per\_s}
legítimamente enorme dominaran la proyección (dos puntos en $\pm 300$, el resto aplastado
en un blob). Recortar al percentil 1--99 y usar \code{RobustScaler} \emph{tampoco} alcanzó:
varias columnas están tan infladas en cero (la mediana y los percentiles 25/75 son 0 para
la mayoría de los flags/ratios) que su IQR es 0, y el \emph{fallback} de \code{RobustScaler}
ante IQR=0 es dejar la columna \textbf{sin escalar} --- su magnitud cruda (microsegundos, en
algunas columnas temporales) volvió a dominar todo. La solución fue \textbf{log1p con
signo} ($\mathrm{sign}(x)\cdot\log(1+|x|)$, necesario porque \code{init\_fwd\_win\_byts}/
\code{init\_bwd\_win\_byts} usan $-1$ como centinela de "desconocido") antes de
\code{StandardScaler}. El preprocesamiento de datos de red de cola pesada no es un detalle
menor de la tubería --- fue, en los hechos, la mitad del trabajo de esta sección.
\end{cbox}

$k=8$ (silhouette 0.531). El 99.99\% del tráfico de ataque de la muestra (22.083 de 22.086
filas) cae en sólo \textbf{2 de los 8 clusters}:

\begin{center}
\begin{tabular}{@{}rrrr@{}}
\toprule
\textbf{Cluster} & \textbf{benigno} & \textbf{ataque} & \textbf{\% ataque} \\
\midrule
2 & 75.051 & 10.974 & 12.8\% \\
5 & 21.400 & 11.109 & \textbf{34.2\%} (vs.\ 7.2\% base) \\
0, 1, 3, 4, 6, 7 & 187.184 & 3 & $\approx$0\% \\
\bottomrule
\end{tabular}
\end{center}

\begin{figure}[ht]
\centering
\includegraphics[width=0.95\textwidth]{../../results/figures/real_data_cluster_pca.png}
\caption{Proyección PCA (2 componentes, 65.4\% de varianza explicada) de la muestra,
coloreada por cluster no supervisado (izquierda) y por etiqueta real (derecha, rojo =
ataque).}
\end{figure}

Ni "sin estructura" ni "separación limpia": el ataque \textbf{se concentra}, no se
distribuye uniformemente entre los 8 clusters, pero tampoco forma un cluster propio --- los
dos clusters que lo contienen son mayoritariamente benignos (87\% y 66\% respectivamente).
Para un clasificador supervisado esto es una noticia mixta: hay estructura no supervisada
que correlaciona con el ataque (más de lo que el AUC per-flujo invertido de
\S\ref{sec:contraste} sugeriría solo), pero separarla limpiamente todavía exige la
etiqueta, no emerge sola.

% ============================================================
\section{Primer benchmark supervisado real}
\label{sec:benchmark}
% ============================================================

Todo lo medido hasta acá evaluaba \emph{un} feature aislado (\S\ref{sec:contraste}) o
estructura \emph{sin} etiqueta (\S\ref{sec:clustering}). Esta sección corre, por primera
vez contra datos reales, los cinco modelos que \code{src/models/benchmark.py} ya tiene
implementados --- con manejo de desbalanceo (\code{class\_weight="balanced"} /
\code{scale\_pos\_weight}) y SHAP --- \textbf{a nivel de flujo}, sin construir grafos ni
ventanas: las funciones de \code{src/models/benchmark.py} sólo piden columnas de features
más \code{has\_attack}, así que el problema de potencia que bloqueó H2
(\S\ref{sec:phase0}, sólo 63--383 ventanas) \textbf{no aplica acá} --- se usan los 7.9M
flujos completos, con \code{src.models.benchmark.temporal\_split} (70/30 cronológico, no
aleatorio).

\begin{cbox}{warn}{Por qué esta tabla no encabeza con ROC-AUC}
Bajo desbalanceo (7.2\%--35.9\% ataque según archivo/split), ROC-AUC puede verse fuerte
aunque el modelo sea mediocre detectando la clase que importa --- pondera los verdaderos
negativos contra una mayoría benigna enorme. La tabla encabeza con \textbf{precisión y
recall de la clase ataque específicamente} (no promediados con la clase benigna, que es
casi siempre más fácil) y PR-AUC, cuya línea base ya es la tasa de positivos del
conjunto de prueba, no 0.5. \code{f1\_macro} (guardado en
\code{results/metrics/real\_data\_benchmark.json} para quien lo quiera) queda deliberadamente
fuera de la tabla por la misma razón: en Infiltración, \code{f1\_macro}=0.687 para
RandomForest suena mejor que el F1(ataque)=0.591 real --- la diferencia es exactamente el
desbalanceo que \code{f1\_macro} esconde.
\end{cbox}

\begin{center}
\begin{tabular}{@{}lrrrrr@{}}
\toprule
\textbf{Modelo} & \textbf{PR-AUC} & \textbf{Precisión(atq)} & \textbf{Recall(atq)} & \textbf{F1(atq)} & \textbf{Tiempo} \\
\midrule
\multicolumn{6}{@{}l}{\textit{Infiltración (231.718 train / 99.309 test, 35.9\% ataque en test)}} \\
RandomForest    & 0.713 & 0.612 & 0.571 & 0.591 & 144.2 s \\
XGBoost         & 0.705 & 0.562 & 0.605 & 0.583 & 2.1 s \\
IsolationForest & 0.317 & 0.245 & 0.056 & 0.091 & 2.4 s \\
Autoencoder     & 0.328 & 0.223 & 0.056 & 0.089 & 4.9 s \\
OneClassSVM$^\dagger$ & 0.332 & 0.357 & 0.106 & 0.163 & 70.4 s \\
\multicolumn{6}{@{}l}{\textit{DDoS-LOIC-HTTP (5.564.122 train / 2.384.624 test, 17.2\% ataque en test)}} \\
RandomForest    & \textbf{1.000} & \textbf{1.000} & \textbf{1.000} & \textbf{1.000} & 341.3 s \\
XGBoost         & \textbf{1.000} & \textbf{1.000} & \textbf{1.000} & \textbf{1.000} & 35.5 s \\
IsolationForest & 0.242 & 0.015 & 0.007 & 0.010 & 74.4 s \\
Autoencoder     & 0.172 & 0.015 & 0.007 & 0.010 & 96.6 s \\
OneClassSVM$^\dagger$ & 0.161 & 0.015 & 0.007 & 0.010 & 1397.0 s \\
\bottomrule
\end{tabular}
\end{center}

\noindent $^\dagger$ Entrenado sobre una muestra sistemática de 50.000 filas benignas, no
el conjunto completo --- costo $O(n^2)$--$O(n^3)$ documentado en el propio
\code{run\_one\_class\_svm} (confirmado: 1397\,s incluso con esa muestra acotada).

\noindent \textit{Nota de procedencia}: Precisión/F1(ataque) de los tres modelos no
supervisados en DDoS se derivan algebraicamente ---
$\text{precisión}=\mathrm{TP}/(\mathrm{TP}+\mathrm{FP})$, con
$\mathrm{TP}=\text{recall}\times n_{\text{ataque}}$ y $\mathrm{FP}=\text{FPR}\times
n_{\text{benigno}}$ --- a partir de recall/FPR y el tamaño exacto de cada clase en el
conjunto de prueba (410.301 ataque / 1.974.323 benigno), ya medidos en una corrida
completa anterior; no son una estimación. RandomForest/XGBoost sí se re-corrieron de
punta a punta (XGBoost, 122\,s con carga de datos incluida) para confirmar los números y
generar la curva real de abajo, después de que el proceso en segundo plano de esta
sesión se interrumpiera externamente tres veces seguidas sobre la corrida completa de
los cinco modelos.

\begin{figure}[ht]
\centering
\includegraphics[width=0.95\textwidth]{../../results/figures/real_data_benchmark_Thursday-01-03-2018_TrafficForML_CICFlowMeter_RandomForest_pr_roc.png}
\caption{Curvas PR y ROC reales, RandomForest sobre Infiltración. La curva PR --- no el
escalar PR-AUC --- muestra dónde importa: precisión $\approx$0.95--0.99 hasta recall
$\approx$0.35, y cae con fuerza después. Al umbral por defecto (0.5) el modelo opera en la
zona de caída (precisión 0.612, recall 0.571) --- ajustar el umbral es un trade-off real,
visible acá y ausente en un solo número de AUC.}
\end{figure}

\begin{figure}[ht]
\centering
\includegraphics[width=0.95\textwidth]{../../results/figures/real_data_benchmark_Thuesday-20-02-2018_TrafficForML_CICFlowMeter_XGBoost_pr_roc.png}
\caption{Curvas PR y ROC reales, XGBoost sobre DDoS. A diferencia de Infiltración, acá la
curva es literalmente el escalón perfecto: precisión = 1.0 para todo recall hasta 1.0.
No es un artefacto de la escala del gráfico --- es el resultado exacto, verificado
corriendo el modelo, no asumido a partir del AUC.}
\end{figure}

\begin{cbox}{crit}{AUC = 1.000 exige el mismo escrutinio que cualquier otro 1.0000 en este proyecto}
\code{CLAUDE.md} documenta, como mistake \#1, que \textbf{este mismo proyecto} ya publicó
un AUC de 1.0000 que resultó ser un confound (\code{n\_nodes} separando corridas, no
ataques). Un 1.000 perfecto en 2.384.624 filas de prueba --- 0 falsos positivos, recall
1.000, \textbf{y también precisión 1.000 y F1(ataque) 1.000}, no sólo el AUC agregado ---
exige la misma sospecha, no una celebración automática. Lo que se verificó antes de
escribir esto:

\begin{itemize}[leftmargin=1.4em,itemsep=1pt]
\item \textbf{No es \code{dst\_port} solo.} El 100\% de los flujos de ataque van al
  puerto 80 --- pero también el 11.8\% de los flujos \emph{benignos} de este mismo
  archivo (868.974 de 7.372.555). Un solo feature categórico no alcanza para separar
  perfectamente.
\item \textbf{La importancia de features está repartida}, no concentrada en una: el top
  20 de RandomForest (\code{results/figures/real\_data\_benchmark\_*\_RandomForest\_importance.png})
  reparte el peso entre \code{totlen\_fwd\_pkts} (0.15), \code{fwd\_pkt\_len\_max} (0.11),
  \code{flow\_iat\_max} (0.09), \code{dst\_port} (0.07), y una decena más --- no hay una
  única columna filtrando la etiqueta.
\item \textbf{La diferencia de volumen es real y enorme}: \code{totlen\_fwd\_pkts} medio
  en ataque = 4.956 bytes vs.\ 382 en benigno (13$\times$) --- LOIC-HTTP mantiene flujos
  de flood de larga duración (hasta 120s) con miles de peticiones repetidas, muy distinto
  de un flujo web normal.
\item \textbf{El contraste con lo no supervisado es aún más marcado que en Infiltración}:
  IsolationForest/Autoencoder/OneClassSVM tienen precisión(ataque) $\approx$1.5\% en
  DDoS --- de cada 100 flujos que marcan como ataque, $\approx$1--2 lo son de verdad. Si
  el 1.000 de RandomForest/XGBoost fuera un artefacto trivial del desbalanceo, se
  esperaría que los no supervisados también se beneficiaran de alguna forma; en cambio,
  sin la etiqueta, el problema es casi tan difícil como encontrar una aguja al azar.
\end{itemize}

\textbf{Conclusión, con la misma honestidad que el resto del documento}: no se encontró
un confound de una sola columna, pero tampoco se agotaron todas las formas posibles de
confound (no se verificó, por ejemplo, superposición de IPs entre train y test). Lo que
sí se puede afirmar es que este resultado es \textbf{consistente, no contradictorio}, con
el argumento central del proyecto (\code{docs/metodologia/REPLANTEO-OPERACIONAL.md}): las
features completas de CICIDS presuponen visibilidad total de paquetes, y bajo ese régimen
limpio la detección \textbf{satura trivialmente} --- ya se había medido así en el
simulador (AUC 1.0000 en fuerza bruta, techo en PR-AUC); esta es la primera vez que se
mide así, con el mismo patrón, en \textbf{datos reales}. No es evidencia nueva a favor del
modelo --- es evidencia de que el régimen limpio no deja margen para medir nada, exactamente
lo que motivó todo el replanteamiento operacional.
\end{cbox}

La disparidad train/test es aún más marcada que en Infiltración: 2.98\% de ataque en
entrenamiento vs.\ 17.2\% en prueba --- el ataque real se concentra hacia el final
cronológico de la captura. Con 5.56M filas de entrenamiento, 2.98\% siguen siendo
$\approx$166.000 flujos de ataque para entrenar, no una muestra degenerada.

\begin{cbox}{good}{En Infiltración: supervisado gana, pero lejos de perfecto}
RandomForest atrapa el 57\% de los ataques reales (recall) y, de lo que marca como
ataque, acierta el 61\% (precisión) --- F1(ataque)=0.591. Los tres baselines no
supervisados, entrenados sólo con tráfico benigno, ni siquiera se acercan: recall
5.6--10.6\%, F1(ataque) 0.09--0.16. Es la misma historia que el clustering de
\S\ref{sec:clustering} ya insinuaba (el ataque no se aísla solo) llevada a su conclusión
lógica: sin la etiqueta, ningún método no supervisado probado encuentra al ataque;
combinando las 48 features \emph{con} la etiqueta, un modelo de árboles sí recupera una
detección real --- lejos de perfecta (recall 57\%, no 100\%), pero muy por encima del AUC
de 0.537 que el feature aislado daba solo, y del recall $\approx$0 de intentarlo sin
etiqueta. La combinación de features hace el trabajo que ningún feature ni ningún cluster
hacía por separado. El patrón de "no supervisado no encuentra nada, supervisado sí" se
repite en DDoS (arriba) --- con una magnitud muy distinta.
\end{cbox}

Nota metodológica honesta: el split temporal (no aleatorio, para evitar fuga de
información) muestra tasas de ataque distintas entre entrenamiento (24.8\%) y prueba
(35.9\%) en el archivo de Infiltración --- el ataque no está distribuido uniformemente en
el tiempo dentro del propio archivo. Es exactamente el tipo de cosa que un split aleatorio
escondería y que este split, más honesto, expone.

% ============================================================
\section{H2 bajo restricción, con datos reales: el experimento preregistrado, por fin corrido}
\label{sec:h2}
% ============================================================

\code{experiments/04\_restriction.py} es el experimento preregistrado de H2
(\code{docs/metodologia/REPLANTEO-OPERACIONAL.md}): ¿la topología degrada \emph{menos} que la
asimetría per-flujo bajo restricción de información (muestreo 1:N, truncamiento temporal)? Su
propio docstring decía, textual, que existe para correr \textit{``sin cambios en cuanto lleguen
datos reales''}. Nunca había corrido contra datos reales --- sólo contra el simulador, donde el
baseline per-flujo satura en 1.0000 y el gate de potencia (\code{MIN\_WINDOWS\_PER\_CLASS=250}) no
se cumple nunca. Esta sección lo corre por primera vez contra \code{Thuesday-20-02-2018}, sin
modificar el gate, el criterio de decisión, ni los controles --- sólo se agregó una fuente de
datos real (\code{--data-path}) y un parámetro de tamaño de ventana (\code{--window-seconds}), con
el mismo split temporal 70/30 (\code{temporal\_split}) ya usado en \S\ref{sec:benchmark}.

\begin{center}
\begin{tabular}{@{}lrrrrl@{}}
\toprule
\textbf{Eje} & \textbf{ventana} & \textbf{n+/n- (min)} & \textbf{per-flujo} & \textbf{topología} & \textbf{gate} \\
\midrule
Muestreo    & 60s & 45 & 0.891--1.000 & 0.514--0.863 & potencia + saturación \\
Truncamiento & 60s & 45 & 1.000 (las 4) & 0.655--0.732 & potencia + saturación \\
Muestreo    & 10s & 2--270 & 0.500--1.000 & 0.500--0.946 & saturación (potencia \emph{sí} cumple en 3/4) \\
Truncamiento & 10s & 270 & 0.995--1.000 & 0.603--0.819 & saturación \\
\bottomrule
\end{tabular}
\end{center}

\begin{cbox}{good}{El hallazgo real: a 10s, el gate de potencia por fin se cumple --- y el de saturación sigue sin ceder}
Con ventanas de 60s (el default de \code{window\_features}), la clase minoritaria del split de
\emph{prueba} queda en 45 ventanas en las 8 corridas --- muy por debajo de 250. Con ventanas de
10s, sube a \textbf{270} en 3 de 4 niveles de muestreo y en los 4 de truncamiento ---
\textbf{la primera vez en la historia de este proyecto que el criterio de $\geq$250 ventanas por
clase se cumple en una corrida real del experimento preregistrado}, no sólo en un barrido
descriptivo (\S\ref{sec:phase0}). Pero el gate de saturación no cede: el baseline per-flujo mide
entre 0.977 y 1.000 en absolutamente todas las combinaciones de muestreo/truncamiento con ventanas
de 10s, salvo en el punto más extremo (muestreo 1:1000), donde los datos colapsan a 2 ventanas de
ataque y todo cae a 0.500 por falta de datos, no por restricción real.
\end{cbox}

\textbf{El patrón direccional, reportado sin certificar un veredicto} (el gate sigue bloqueando
uno formal, y esta sesión no lo fuerza con \code{--force} salvo para depuración): en las
\textbf{16 combinaciones} corridas, la topología queda \emph{por debajo} del per-flujo en las 16
--- nunca al revés, ni siquiera en el punto de mayor restricción real (muestreo 1:100, ventana
10s: topología 0.946 vs.\ per-flujo 0.977). Si el gate hubiera cumplido, la lectura honesta según
la propia regla de decisión del script sería \code{H2 REJECTED} (delta $<$ 0.05 en todos los
niveles) --- no \code{H2 SUPPORTED}. Esto es lo opuesto de lo que H2 predice (que la topología
degrade \emph{menos}); en este archivo específico, con este tipo de ataque específico
(DDoS-LOIC-HTTP), el per-flujo resulta más robusto a la restricción de lo que la teoría
anticipaba, no menos.

\begin{cbox}{warn}{Por qué esto no cierra la pregunta}
Un solo archivo, un solo tipo de ataque, ningún control C2/C3 fallado pero un gate de saturación
que nunca cede: no alcanza para generalizar. Lo que sí cambia respecto a antes de esta sesión es
que el bloqueo pasó de ser \textit{``no hay suficientes ventanas para intentarlo''} a
\textit{``se puede intentar, y el per-flujo sigue sin dejar margen que medir''} --- un diagnóstico
más específico y más cerca de una respuesta real. El siguiente paso honesto es un segundo día con
IPs (opciones (b)/(c) de \code{SETUP-LOCAL.md}) o un tipo de ataque distinto al de LOIC-HTTP, para
saber si el patrón de saturación es del archivo o del fenómeno.
\end{cbox}

\subsection{Re-corrida con windowing canónico: ¿cambia el instrumento el resultado?}
\label{sec:h2-canonical}

\code{scripts/run\_h2\_real.py} (nuevo) repite exactamente los mismos 8 niveles (4 de muestreo, 4
de truncamiento) y las mismas dos ventanas (60s, 10s) de la corrida anterior, pero usando
\code{src.audit.canonical} (el windowing corregido por P0.5) en vez de
\code{src.features.windows.window\_features}, que re-derivaba los límites de ventana desde el
\code{ts.min()} de cada frame ya restringido --- exactamente el bug de deriva de ancla que P0.5
corrigió. También agrega el bootstrap 95\,\% CI pareado que \code{experiments/04\_restriction.py}
nunca implementó (\S\ref{sec:h2}, el mismo método que
\code{docs/methodology/H2\_SYNTHETIC\_FINDINGS.md} introdujo para el resultado sintético). El split
train/test sigue siendo temporal (70/30) pero ahora alineado a un borde de ventana canónica, no a
un corte de fila cruda, para que ninguna ventana quede partida entre los dos conjuntos.

\textbf{Nota sobre las etiquetas de truncamiento a \code{window\_seconds=10}.} Los niveles de
truncamiento se heredan sin cambios de \code{experiments/04\_restriction.py}: son
\emph{fracciones} (1/60, 5/60, 10/60, 1.0), no segundos absolutos, definidas cuando la única
ventana probada era de 60s. A \code{window\_seconds=10} las mismas fracciones corresponden a
$\approx$0.17s/0.83s/1.67s/10s observados, no literalmente 1s/5s/10s/60s --- se preserva la
convención original (y su etiqueta) tal cual para que la comparación con la tabla anterior sea
exacta, pero el valor real detrás de cada etiqueta a 10s es el que aquí se aclara, no el que el
nombre sugiere.

\begin{center}
\small
\begin{tabular}{@{}lrrrrl@{}}
\multicolumn{6}{l}{\textbf{Instrumento antiguo (PR \#24, windowing pre-P0.5)}} \\
\toprule
\textbf{Eje} & \textbf{ventana} & \textbf{n+/n- (min)} & \textbf{per-flujo} & \textbf{topología} & \textbf{gate} \\
\midrule
Muestreo    & 60s & 45 & 0.891--1.000 & 0.514--0.863 & potencia + saturación \\
Truncamiento & 60s & 45 & 1.000 (las 4) & 0.655--0.732 & potencia + saturación \\
Muestreo    & 10s & 2--270 & 0.500--1.000 & 0.500--0.946 & saturación (potencia \emph{sí} cumple en 3/4) \\
Truncamiento & 10s & 270 & 0.995--1.000 & 0.603--0.819 & saturación \\
\bottomrule
\end{tabular}

\vspace{0.6em}

\begin{tabular}{@{}lrrrrl@{}}
\multicolumn{6}{l}{\textbf{Windowing canónico (esta sesión, \code{scripts/run\_h2\_real.py})}} \\
\toprule
\textbf{Eje} & \textbf{nivel} & \textbf{n+/n-} & \textbf{per-flujo} & \textbf{topología} & \textbf{gate} \\
\midrule
Muestreo (60s)    & 1:1--1:1000 & 63/(123--153) & --- (nan) & --- (nan) & potencia \\
Truncamiento (60s) & 1s--60s    & 63/153        & --- (nan) & --- (nan) & potencia \\
Muestreo (10s)    & 1:1   & 378/913 & 0.5174 & 0.4899 & \textbf{OK} \\
Muestreo (10s)    & 1:10  & 378/913 & 0.7891 & 0.4609 & \textbf{OK} \\
Muestreo (10s)    & 1:100 & 378/890 & 0.4571 & 0.5050 & \textbf{OK} \\
Muestreo (10s)    & 1:1000 & 4/0    & --- (nan) & --- (nan) & potencia \\
Truncamiento (10s) & $\approx$0.17s & 378/913 & 0.8082 & 0.4617 & \textbf{OK} \\
Truncamiento (10s) & $\approx$0.83s & 378/913 & 0.8082 & 0.4617 & \textbf{OK} \\
Truncamiento (10s) & $\approx$1.67s & 378/913 & 0.6921 & 0.4619 & \textbf{OK} \\
Truncamiento (10s) & 10s (completo) & 378/913 & 0.5174 & 0.4899 & \textbf{OK} \\
\bottomrule
\end{tabular}
\end{center}

\begin{cbox}{good}{El instrumento sí cambió el resultado --- y de una forma que refuerza, no debilita, la conclusión}
A \code{window\_seconds=60} nada cambió: el gate de potencia sigue bloqueando todo (63 ventanas de
ataque en el nuevo conteo contra 45 en el viejo --- la cifra exacta se movió porque el ancla ahora
es fija, pero la conclusión no). A \code{window\_seconds=10} el cambio es grande: el baseline
per-flujo \textbf{ya no satura en absoluto}. Bajo el windowing viejo medía 0.977--1.000 en casi
todas las combinaciones; bajo el windowing corregido mide entre \textbf{0.4571 y 0.8082}, y en
información completa (1:1 muestreo $\equiv$ 10s truncamiento, ambas filas dan exactamente los
mismos números: 0.5174 per-flujo, 0.4899 topología --- la verificación interna de que ambos caminos
de código coinciden en el caso sin restricción) el per-flujo está prácticamente en el azar. Esto
implica que el bug de deriva de ancla que P0.5 corrigió no era un artefacto menor: en este archivo,
\textbf{estaba inflando artificialmente la separabilidad per-flujo hasta casi 1.0}. Y por primera
vez en la historia de este proyecto, \textbf{7 de las 8 combinaciones a 10s despejan ambos gates a
la vez} (potencia $\geq$250 y saturación $\leq$0.90) --- sólo el muestreo 1:1000 sigue bloqueado
(la clase benigna colapsa a 0 ventanas).
\end{cbox}

\textbf{Veredicto certificado, por primera vez con datos reales:} en las 7 combinaciones que
despejan ambos gates, la regla de decisión pre-registrada da \code{H2 REJECTED} en las 7 ---
\texttt{delta}(topología $-$ per-flujo) nunca alcanza $+0.05$. En 6 de las 7, \texttt{delta} es
claramente negativo (per-flujo gana), con el CI bootstrap de 2000 remuestreos completamente por
debajo de cero. La excepción interesante es \textbf{muestreo 1:100}: \texttt{delta}=$+0.0479$,
CI=$[+0.0221, +0.0736]$ --- el intervalo \emph{excluye} cero (topología gana de forma
estadísticamente real, no es ruido), pero el valor no llega al umbral de $0.05$ que la regla
pre-registrada exige para contar como evidencia a favor de H2. Es el punto más cercano que la
topología ha estado nunca de superar al per-flujo, en cualquiera de los dos instrumentos
(sintético o real) --- y aun así, no alcanza.

\begin{cbox}{warn}{Exploratorio, no parte de la regla de decisión de H2: ¿ayuda combinar features?}
Igual que en el resultado sintético (\code{H2\_SYNTHETIC\_FINDINGS.md}), se calculó el mismo CI
bootstrap pareado para \texttt{combined} $-$ \texttt{per-flujo}, exclusivamente post-hoc. A
diferencia del caso sintético (donde el CI incluía cero en los 5 niveles), aquí el patrón es más
fuerte: muestreo 1:10 da \texttt{delta}=$+0.0942$, CI=$[+0.0638, +0.1249]$ (excluye cero, y por
encima del umbral de 0.05 que H2 exige --- pero esta comparación nunca fue la hipótesis
pre-registrada); muestreo 1:100 da \texttt{delta}=$+0.0795$, CI=$[+0.0508, +0.1073]$; truncamiento
$\approx$1.67s da \texttt{delta}=$+0.0756$, CI=$[+0.0524, +0.0982]$. Las dos filas de
truncamiento más cortas ($\approx$0.17s/0.83s) dan \texttt{delta}=$-0.0281$, CI=$[-0.0647,+0.0098]$
(incluye cero). Que esta señal exploratoria sea más fuerte en datos reales que en el sintético es,
en sí, un hallazgo a anotar --- no a convertir retroactivamente en la pregunta confirmatoria.
\end{cbox}

% ============================================================
\section{Completitud: las 10 capturas}
\label{sec:completitud}
% ============================================================

Todo lo anterior en este documento --- el AUC invertido, el grafo, el clustering, el benchmark
supervisado, H2 --- se midió sobre 2 de los 10 archivos del bucket. Esta sección corre los mismos
dos scripts (\code{analyze\_real\_data.py}, \code{benchmark\_real\_data.py} --- \textbf{sin
modificar ninguno}, ya aceptaban un directorio completo) sobre los 10, para responder la pregunta
obvia: ¿el patrón se repite en los otros 8 tipos de ataque, o es específico de DDoS-LOIC-HTTP e
Infiltración?

Primero, una corrección de alcance: la documentación de este repositorio (\code{PLAN-DE-RECURSOS.md},
\code{SETUP-LOCAL.md}) asumía \textit{``\textasciitilde 6.5 GB comprimido / \textasciitilde 60 GB
descomprimido''} para el dataset completo --- un supuesto nunca verificado. Sumando los 10 archivos
reales del bucket (\code{aws s3 ls}, sin paso de descompresión: son CSV planos): \textbf{6.886.649.507
bytes $\approx$ 6.4 GiB en total}. Entran los 10 a la vez en cualquier máquina con SSD moderno;
no hace falta procesar por lotes como la documentación anterior recomendaba.

\begin{center}
\begin{longtable}{@{}lp{4.4cm}rrr@{}}
\toprule
\textbf{Archivo} & \textbf{Etiqueta(s) real(es)} & \textbf{\% ataque} & \textbf{AUC per-flujo} & \\
\midrule
\endhead
Friday-02-03 & Bot & 27.1\% & \textbf{0.838} & \\
Friday-16-02 & DoS-Hulk, DoS-SlowHTTPTest & 50.4\% & \textbf{0.023} & \\
Friday-23-02 & Brute Force-Web/-XSS, SQL Injection & 0.05\% & 0.694 & \\
Thuesday-20-02 & DDoS-LOIC-HTTP & 7.2\% & 0.381 & \\
Thursday-01-03 & Infiltration & 28.1\% & 0.537 & \\
Thursday-15-02 & DoS-GoldenEye, DoS-Slowloris & 4.9\% & 0.590 & \\
Thursday-22-02 & Brute Force-Web/-XSS, SQL Injection & 0.03\% & 0.633 & \\
Wednesday-14-02 & SSH-/FTP-BruteForce & 19.0\% & 0.631 & \\
Wednesday-21-02 & DDoS attack-HOIC, -LOIC-UDP & 65.0\% & \textbf{0.003} & \\
Wednesday-28-02 & Infiltration & 11.3\% & 0.494 & \\
\bottomrule
\end{longtable}
\end{center}

\begin{cbox}{crit}{El patrón no es "la asimetría se invierte" --- es "se invierte según el mecanismo del ataque"}
Con las 10 etiquetas reales confirmadas (no de memoria: \code{load\_csv(f).select('label').collect()},
el mismo comando para las 10), aparece una estructura que 2 archivos no podían mostrar:

\begin{itemize}[leftmargin=1.4em,itemsep=2pt]
\item \textbf{Los ataques de \emph{flood} (DoS/DDoS) invierten la asimetría, y lo hacen todos}: DoS-Hulk
  (0.023), DDoS-HOIC/LOIC-UDP (0.003), DDoS-LOIC-HTTP (0.381) --- las tres peores puntuaciones de
  las 10, y las tres son variantes de saturar un servidor con peticiones. Wednesday-21-02
  (\textbf{0.003}) es más extremo que cualquier número medido antes en este documento: casi
  perfectamente invertido.
\item \textbf{Bot y Brute Force/SQL Injection preservan la dirección esperada}: Bot llega a 0.838 ---
  el mejor AUC per-flujo real medido en todo el dataset, mejor que cualquier cosa reportada hasta
  ahora en este documento. Brute Force/SQLi y SSH/FTP-BruteForce quedan en 0.63--0.69, débiles pero
  en la dirección correcta.
\item \textbf{Infiltration es internamente consistente}: dos capturas distintas del mismo tipo de
  ataque (Thursday-01-03, Wednesday-28-02) dan 0.537 y 0.494 --- ambas cerca del azar, sin la
  inversión de los floods ni la fuerza del Bot. Es la clase de repetibilidad entre corridas
  independientes que un hallazgo real debería mostrar.
\end{itemize}

Lectura mecánica, coherente con lo que cada ataque hace: un flood manda peticiones y \emph{recibe}
la respuesta/error del servidor (fwd $\ll$ bwd, la dirección opuesta al supuesto de diseño); un bot
o un ataque de fuerza bruta/inyección \emph{envía} el payload activamente (fwd $\gg$ bwd, la
dirección que \code{A\_dir} siempre asumió). \textbf{La asimetría per-flujo no es una señal rota ---
es una señal cuya dirección depende del mecanismo del ataque}, algo que 2 archivos (uno de cada
extremo, por casualidad) no podían mostrar y que ninguna de las 10 referencias ya verificadas en
\code{paper/references.bib} documenta explícitamente para este dataset.
\end{cbox}

También consistente en las 10: 49/49 \code{CLASSIC\_FEATURES} presentes, 24/49 direccionales
(estructural, no depende del archivo), 6--9 grupos redundantes al umbral 0.80
(\S\ref{sec:corr} usaba 6 y 9 como los extremos de 2 archivos; con los 10 el rango es el mismo,
6--9 --- el umbral generaliza). Headers duplicados incrustados aparecieron en \textbf{3 de 10}
archivos (Thursday-01-03, Friday-16-02, Wednesday-28-02), no sólo en el que motivó el fix original
de \code{load\_csv()} (PR \#17) --- confirma que el bug que ese fix cerró no era un caso aislado.

El mismo \code{benchmark\_real\_data.py} de \S\ref{sec:benchmark} (RandomForest, XGBoost,
IsolationForest, Autoencoder, OneClassSVM, sin cambiar ninguno) corrido sobre los 10, mismo
90/10 split temporal por archivo, reporta P/R/F1(ataque) reales de cada uno --- no PR-AUC solo:

\begin{center}
\begin{longtable}{@{}lp{4.0cm}lrrr@{}}
\toprule
\textbf{Archivo} & \textbf{Ataque} & \textbf{Mejor modelo} & \textbf{P(ataque)} & \textbf{R(ataque)} & \textbf{F1(ataque)} \\
\midrule
\endhead
Thursday-01-03 & Infiltration & RandomForest & 0.612 & 0.571 & 0.591 \\
Friday-02-03 & Bot & XGBoost & 1.000 & 1.000 & \textbf{1.000} \\
Friday-16-02 & DoS-Hulk, DoS-SlowHTTPTest & OneClassSVM & 0.948 & 1.000 & \textbf{0.973} \\
Friday-23-02 & Brute Force-Web/-XSS, SQL Injection & XGBoost & 0.395 & 0.842 & 0.538 \\
Thuesday-20-02 & DDoS-LOIC-HTTP & XGBoost & 1.000 & 1.000 & \textbf{1.000} \\
Thursday-15-02 & DoS-GoldenEye, DoS-Slowloris & XGBoost & 1.000 & 0.761 & 0.864 \\
Thursday-22-02 & Brute Force-Web/-XSS, SQL Injection & XGBoost & 0.969 & 0.992 & \textbf{0.981} \\
Wednesday-14-02 & SSH-/FTP-BruteForce & RandomForest & 1.000 & 1.000 & \textbf{1.000} \\
Wednesday-21-02 & DDoS attack-HOIC, -LOIC-UDP & RandomForest & 1.000 & 0.988 & \textbf{0.994} \\
Wednesday-28-02 & Infiltration & XGBoost & 0.276 & 0.397 & 0.326 \\
\bottomrule
\end{longtable}
\end{center}

\begin{cbox}{crit}{La saturación bajo visibilidad completa generaliza --- con dos excepciones reales, no una}
\S\ref{sec:benchmark} ya había mostrado F1(ataque) = 1.000 en el archivo de DDoS bajo visibilidad
completa (todas las columnas de CICFlowMeter, no sólo la asimetría per-flujo) y se preguntaba si
era el archivo o el fenómeno. Con los 10, la respuesta es clara: \textbf{7 de 10 archivos saturan}
(F1(ataque) $\geq$ 0.86) bajo un clasificador de flujo completo, cubriendo tipos de ataque tan
distintos como Bot, DoS-Hulk/SlowHTTPTest, DDoS-LOIC-HTTP, DoS-GoldenEye/Slowloris, SSH/FTP-BruteForce
y DDoS-HOIC/LOIC-UDP --- la saturación no es un artefacto de los 2 archivos originales, es el
régimen dominante de CICIDS-2018 a nivel de flujo.

Las excepciones son las que importan:
\begin{itemize}[leftmargin=1.4em,itemsep=2pt]
\item \textbf{Infiltration falla en las dos capturas independientes}: Thursday-01-03 (F1=0.591) y
  Wednesday-28-02 (F1=0.326) --- consistente con el AUC per-flujo cercano al azar de la tabla
  anterior, y consistente entre sí (la misma repetibilidad que ya se destacó ahí). Es el único tipo
  de ataque que un clasificador de flujo con visibilidad completa no resuelve bien en ninguna de
  sus dos apariciones.
\item \textbf{Dos archivos con la misma etiqueta (Brute Force-Web/-XSS, SQL Injection) y tasas de
  ataque casi idénticas dan resultados opuestos}: Thursday-22-02 satura (F1=0.981, 128 positivos en
  test) pero Friday-23-02 no (F1=0.538, 177 positivos en test) --- conteos de la clase positiva del
  mismo orden de magnitud, así que no es sólo un problema de tamaño de muestra. Esto queda como
  pregunta abierta, no como hallazgo resuelto: la etiqueta agrupa sub-ataques de \code{CICFlowMeter}
  (Brute Force-Web, Brute Force-XSS, SQL Injection son tres etiquetas de ataque distintas colapsadas
  en una sola categoría por esta tabla) que pueden no ser comparables entre capturas, y el propio
  \code{scripts/benchmark\_real\_data.py} no las separa. Documentado aquí en vez de promediado
  silenciosamente en la fila de la tabla.
\end{itemize}
\end{cbox}

% ============================================================
\section{Síntesis: desbalanceo, asimetría, features, y estado del arte}
\label{sec:sintesis}
% ============================================================

Los hallazgos de este documento no viven por separado:

\begin{itemize}[leftmargin=1.4em]
\item \textbf{Desbalanceo.} 7.2\% de tasa de ataque en la muestra completa (DDoS) --- el
  régimen para el que Chawla et al.\ 2002 (SMOTE, \code{chawla2002smote} en
  \code{paper/references.bib}, ya verificado) es la referencia estándar del área. El clustering
  (\S\ref{sec:clustering}) muestra que el desbalanceo no es sólo un problema de conteo: ni
  siquiera un algoritmo no supervisado, sin ver la etiqueta, aísla el ataque en su propio
  grupo --- se concentra en 2 de 8, mezclado con mayoría benigna.
\item \textbf{Asimetría.} El AUC per-flujo invertido (0.381, \S\ref{sec:contraste}) tiene
  ahora una ilustración concreta: un flujo de ataque real de esta muestra manda 20 bytes y
  recibe 964 (\code{totlen\_fwd\_pkts}=20, \code{totlen\_bwd\_pkts}=964) --- exactamente la
  dirección opuesta al supuesto "atacante manda más" sobre el que se construyó
  $A_{\text{dir}}$.
\item \textbf{Grafo/topología.} El grafo real (\S\ref{sec:grafo}) muestra el mecanismo
  detrás del confound de C1 en \S\ref{sec:ddos-results}: un componente de ataque aislado
  agrega estructura contable (más aristas, un sub-grafo extra) que no es la asimetría
  relacional que el proyecto busca medir.
\item \textbf{Estado del arte.} Akoglu et al.\ 2015 (grafos para detección de anomalías,
  \code{akoglu2015graph}) y Zhong et al.\ 2024 (GNN para IDS, survey,
  \code{zhong2024gnnids}) --- ambos ya en \code{paper/references.bib}, verificados ---
  describen exactamente esta tensión: la señal
  estructural existe, pero rara vez está limpiamente separada sin aprendizaje supervisado o
  semi-supervisado sobre la estructura, no sólo agregación hecha a mano
  (\code{top\_reciprocity}, \code{top\_indeg\_gini}, etc.).
\item \textbf{Modelos de ML.} Un clasificador real sobre estos datos enfrenta, a la vez:
  49 features con 9 grupos redundantes en este archivo (\S\ref{sec:corr}), una clase minoritaria al 7.2\%,
  una señal de dirección invertida respecto al supuesto de diseño, y una estructura de
  cluster que ayuda pero no resuelve la separación. El benchmark supervisado
  (\S\ref{sec:benchmark}) es la respuesta empírica, medida en precisión/recall/F1 de la
  clase ataque --- no en un ROC-AUC agregado que el desbalanceo puede inflar: en DDoS,
  RandomForest/XGBoost combinando todo llegan a precisión=recall=F1=1.000 --- saturación,
  no margen, exactamente lo que el régimen de visibilidad completa de CICIDS predice; en
  Infiltración, más débil (F1(ataque) 0.58--0.59, recall 57--60\%) pero muy por encima de
  cualquier pieza aislada, y de intentarlo sin etiqueta (F1(ataque) 0.01--0.16 en ambos
  archivos, precisión $\approx$1.5\% en DDoS --- casi tan malo como azar). Ningún hallazgo
  de este documento dice que el problema sea irresoluble bajo visibilidad completa ---
  dice que esa visibilidad no existe en producción, que es exactamente el punto que
  \code{REPLANTEO-OPERACIONAL.md} ya hacía antes de tener datos reales para confirmarlo.
\end{itemize}

% ============================================================
\section{Metodología: cómo se midió sin colapsar la máquina}
\label{sec:metodo}
% ============================================================

El archivo DDoS pesa 3.8\,GB / 7.9M filas. Ninguna función de este análisis materializa el
frame completo de 84 columnas:

\begin{itemize}[leftmargin=1.4em]
\item Conteos, nulos, medias y percentiles: una sola pasada perezosa
  (\code{LazyFrame.select(...).collect()}) con $\sim$9 agregaciones por columna --- \code{polars}
  nunca construye el \code{DataFrame} completo para esto.
\item Correlación: requiere una matriz materializada, así que se toma una muestra
  \textbf{sistemática determinística} (cada $k$-ésima fila, $k = \lfloor n/750000 \rfloor$)
  --- sin semilla, sin RNG: el mismo archivo siempre da la misma muestra.
\item Barrido de ventanas: sólo se materializan las 6 columnas que
  \code{window\_features} realmente toca (\code{timestamp}, \code{src\_ip}, \code{dst\_ip},
  \code{is\_attack}, \code{totlen\_fwd\_pkts}, \code{totlen\_bwd\_pkts}), no las 84.
\item Grafo de interacciones (\S\ref{sec:grafo}): se filtran las 7.9M filas a sólo las dos
  ventanas de 60s elegidas \emph{antes} de llamar a \code{build\_temporal\_graphs} --- nunca
  se construyen las $\sim$720 ventanas completas para dibujar dos.
\item Clustering (\S\ref{sec:clustering}): mismo principio de muestra acotada que la
  correlación, nunca los 48 features $\times$ 7.9M filas completos.
\end{itemize}

Tiempo y memoria observados (no prometidos, medidos con \code{Get-Process} durante la
corrida):

\begin{center}
\begin{tabular}{@{}lrr@{}}
\toprule
& \textbf{Infiltración (331k filas)} & \textbf{DDoS (7.9M filas)} \\
\midrule
Tiempo total (\code{analyze\_real\_data.py}) & 5.7 s & 668.5 s ($\approx$11.1 min) \\
Tiempo total (\code{explore\_real\_data\_structure.py}) & --- & 82.7 s \\
RSS pico observado & \multicolumn{2}{c}{$\approx$8 GB (sobre 32 GB disponibles en esta máquina)} \\
\bottomrule
\end{tabular}
\end{center}

El costo no está en las agregaciones perezosas (segundos) sino en el barrido de
\S\ref{sec:phase0}: \code{window\_features} recorre, en Python puro, cada ventana de cada
una de las 15 combinaciones de la grilla --- son 15 pasadas completas sobre hasta 3.873
ventanas cada una, no una sola. Es el costo esperado de reusar la función de producción
tal cual está escrita (que prioriza claridad sobre velocidad de barrido) en vez de
reimplementar una versión vectorizada sólo para este análisis --- una decisión de
``código limpio'' antes que ``código rápido'', razonable para una corrida puntual de
$\sim$11 minutos, no para un barrido que se fuera a repetir con frecuencia.

% ============================================================
\section{Implicaciones para el estado actual del proyecto}
% ============================================================

\begin{enumerate}[leftmargin=1.4em]
\item El bug de \code{load\_csv()} (\S3) significa que \textbf{ningún resultado publicado
  hasta hoy pudo haber usado \code{Thursday-01-03-2018} correctamente} --- pero como no
  existían datos reales antes de esta sesión, no invalida ningún hallazgo previo; sí es
  un fix necesario antes de que alguien lo use.
\item El AUC per-flujo real (0.537 Infiltración, \textbf{0.381 DDoS}, \S\ref{sec:contraste})
  no sólo confirma que el simulador sobreestima la señal per-flujo --- en DDoS mide una
  \emph{inversión de dirección} que el argumento teórico de
  \code{docs/metodologia/REPLANTEO-OPERACIONAL.md} no predecía, sólo la degradación. Vale
  la pena revisar si \code{A\_dir}/la asimetría per-flujo asumen una dirección de ataque
  que no es universal entre familias de ataque.
\item El único archivo con IPs \textbf{no es un fracaso categórico de los criterios de
  Fase 0}, como parecía con el único punto (60s) que \code{preflight\_data.py} probaba
  --- el barrido (\S\ref{sec:phase0}) muestra que \code{window=10s} sí cumple el criterio
  3, y que el criterio 4 queda cerca (8.8--9.8\% contra 10\%) en todo el rango probado, no
  lejos. Antes de invocar las opciones (b)/(c)/(d) de \code{SETUP-LOCAL.md}, vale la pena
  probar ventanas aún más finas sobre este mismo archivo.
\item El AUC a nivel de ventana de la topología (\S\ref{sec:ddos-results}) no se despega
  del control de tamaño C1 --- no hay, todavía, evidencia limpia de que la topología
  aporte algo que el tamaño de ventana no explique. Repetir esta medición con potencia
  adecuada ($\geq$250 ventanas/clase) es el paso que la falta actual de un segundo día con
  IPs bloquea directamente.
\item El grafo real (\S\ref{sec:grafo}) y el clustering (\S\ref{sec:clustering}) dan, por
  primera vez, un mecanismo concreto para el confound de C1: el componente de ataque es un
  sub-grafo aislado, y el ataque se concentra (no se distribuye uniforme, pero tampoco se
  aísla del todo) en 2 de 8 clusters no supervisados. Ninguna de las dos cosas se sabía
  antes de esta sesión pese a ser, en el caso del grafo, el objeto que dio nombre al
  proyecto.
\item El benchmark supervisado (\S\ref{sec:benchmark}) es el primer resultado real de
  ``qué tan bien se puede detectar esto'', medido donde importa --- precisión y recall de
  la clase ataque, no un ROC-AUC que el desbalanceo puede inflar. La respuesta depende
  enteramente del régimen: bajo visibilidad completa (lo único que hay hoy), DDoS satura
  (precisión=recall=F1=1.000, verificado que no es un confound de una sola columna) e
  Infiltración queda fuerte pero no perfecta (F1(ataque) 0.58--0.59).
\item El experimento preregistrado H2 (\S\ref{sec:h2}, \S\ref{sec:h2-canonical}) ya tiene
  veredicto certificado con datos reales, no sólo un patrón crudo. La primera corrida (PR
  \#24, windowing pre-P0.5) quedó bloqueada por saturación del per-flujo (0.977--1.000) en
  las 16 combinaciones con potencia suficiente. Repetida con el windowing corregido por
  P0.5 (\S\ref{sec:h2-canonical}), el per-flujo \textbf{deja de saturar por completo}
  (0.4571--0.8082) --- el bug de deriva de ancla estaba inflando artificialmente la
  separabilidad --- y 7 de las 8 combinaciones a 10s despejan ambos gates a la vez. La
  regla de decisión pre-registrada da \code{H2 REJECTED} en las 7. Esto ya no es un
  diagnóstico más específico: es la respuesta que \S\ref{sec:h2} decía que faltaba,
  obtenida sobre el mismo archivo y el mismo tipo de ataque que ya se tenían, cambiando
  únicamente el instrumento de medición. La síntesis completa, junto con el resultado
  sintético equivalente, está en la \S\ref{sec:conclusion}.
\end{enumerate}

% ============================================================
\section{Conclusión}
\label{sec:conclusion}
% ============================================================

\textbf{H2 --- ¿la topología degrada menos que la asimetría direccional per-flujo bajo
restricción de información, y la brecha se amplía?--- queda RECHAZADA}, con evidencia
convergente de dos instrumentos construidos de forma independiente.

\textbf{(i) Generador sintético v2, prueba final con potencia adecuada}
(\code{docs/methodology/H2\_SYNTHETIC\_FINDINGS.md}): en información completa la topología
pierde claramente contra el per-flujo (\texttt{delta}=$-0.2067$,
CI=$[-0.2648,-0.1505]$). Bajo restricción la brecha se estrecha hacia la paridad, nunca
hacia favorecer a la topología, en ningún nivel dentro del alcance evaluado.

\textbf{(ii) Un archivo real de CICIDS-2018 (DDoS-LOIC-HTTP), re-corrida con windowing
canónico esta sesión} (\S\ref{sec:h2-canonical}): veredicto certificado por primera vez ---
\code{H2 REJECTED} en las 7 de 8 combinaciones a \code{window\_seconds=10} que despejan
ambos gates (potencia y saturación). La excepción es muestreo 1:100, donde la topología
gana con significancia estadística real (CI excluye cero) pero por debajo del umbral de
$0.05$ que la regla pre-registrada exige.

Por la propia regla pre-registrada de este proyecto (\code{CLAUDE.md}): \textit{``If
topology degrades equally or worse, the framing does not get amended again --- the project
concludes and the honest result is per-flow only.''} El proyecto concluye. La asimetría
direccional per-flujo es el hallazgo honesto; la topología queda retenida como una
hipótesis auditada y rechazada (la misma postura que \code{CLAUDE.md} ya le da a
\code{src/iag/}: ``kept as an audited baseline, not a proposal'', ahora extendida al nivel
topológico en general, no sólo al índice IAG retirado).

Esta conclusión se limita a lo efectivamente medido: el tráfico sintético del Generador v2
bajo los valores candidatos congelados, y una captura real de un solo tipo de ataque
(DDoS-LOIC-HTTP). No se reclama que generalice a familias de ataque no evaluadas.

\end{document}
